\documentclass[11pt]{amsart}

\usepackage[T1]{fontenc}
\usepackage[utf8]{inputenc}
\usepackage{lmodern}
\usepackage{amsmath}
\usepackage{amssymb}
\usepackage{amsfonts}
\usepackage{amsthm}
\usepackage{mathtools}
\usepackage{geometry}
\usepackage{enumitem}
\usepackage{booktabs}
\geometry{margin=1in}
\usepackage[hidelinks]{hyperref}

\numberwithin{equation}{section}

\theoremstyle{plain}
\newtheorem{theorem}{Theorem}[section]
\newtheorem{proposition}[theorem]{Proposition}
\newtheorem{lemma}[theorem]{Lemma}
\newtheorem{corollary}[theorem]{Corollary}

\theoremstyle{definition}
\newtheorem{definition}[theorem]{Definition}

\theoremstyle{remark}
\newtheorem{remark}[theorem]{Remark}

\newcommand{\dd}{\,\mathrm d}
\newcommand{\RH}{\mathrm{RH}}

\title[A universal $-\beta$ law]{A universal $-\beta$ law for discrete Hankel determinants,\\
with the full leading asymptotics\\
of Hankel determinants of zeta values}

\author{}
\date{Version 05}

\begin{document}

\begin{abstract}
Discrete Hankel determinants built from Dirichlet series decay on the scale $N^{2}\log N$.
That scale has been known since Monien's Coulomb-gas analysis of $\det(\zeta(i+j+r))$, and
Haynes and Zudilin recorded the \emph{expected} leading law
$\log\det(\zeta(a(i+j)+b))=-a\,n^{2}\log n+O(n^{2})$ while proving, by an elementary argument, the
one-sided estimate with coefficient $-a/2$.

This paper asks a different question of the same object: not how small the determinant is, but
\emph{what its leading constant is a function of}.

We prove that for \emph{any} positive discrete measure whose support approaches its single
accumulation endpoint at a regularly varying rate
$1-x(t)=t^{-\beta}L(t)$, and which satisfies only two coarse combinatorial conditions on the
node scale and on the existence of a separated subfamily, the monomial Hankel determinant obeys
\[
\lim_{N\to\infty}\frac{\log D_{N}}{N^{2}\log N}=-\beta .
\]
The leading constant is therefore a single geometric invariant, the \emph{endpoint exponent},
and nothing else: no counting law, no arithmetic, no lattice structure and no weight
information survives at that order. Both bounds are elementary, and the lower bound is a single
Cauchy--Binet term made legitimate by positivity.

Two evaluations follow. First, unconditionally,
$\log\det\bigl(\zeta(a(i+j)+b)\bigr)_{1\le i,j\le n}=-a\,n^{2}\log n+O(n^{2})$, the expected law of
\eqref{eq:HZexpected} in full and in two-sided form, which also doubles the quantitative Diophantine
constant obtainable by that route, from $\liminf\log q_{n}/(n\log n)\ge a/4$ to $\ge a/2$.
Second, the Li--Cayley
map $x(t)=(4t^{2}-1)/(4t^{2}+1)$, which carries the ordinates of the nontrivial zeros of $\zeta$
to $(-1,1)$, has endpoint exponent $\beta=2$; hence every positive Gram system supported on those
Cayley images has determinant constant exactly $-2$, and the Riemann--von Mangoldt refinement
$T\log T$ cannot appear before the second order. The value $\beta=2$ is a one-line evaluation of a
theorem whose content lies entirely in the map $\beta\mapsto-\beta$.

A numerical section confirms the constant to three digits and, more importantly, documents a
trap: naive local exponent estimators are biased at order $\log\log N/\log N$ by the slowly
varying factor, so finite-$N$ fits drift far from $-\beta$ even when the theorem is exact.
\end{abstract}

\maketitle

\section{Introduction}

\subsection{The scale is old; the constant is not}

Let $f(s)=\sum_{m\ge1}a_{m}m^{-s}$ be a Dirichlet series and let
\[
H_{n}^{(a,b)}[f]:=\det_{1\le i,j\le n}\bigl(f(a(i+j)+b)\bigr).
\]
Monien \cite{Monien2009} derived a discrete Vandermonde representation for such determinants and,
by a Coulomb-gas variational analysis, obtained the leading scale
\begin{equation}
\log H_{n}^{(r)}[\zeta]=-n^{2}\Bigl(\log (2n)-\tfrac32\Bigr)+o(n^{2}),
\label{eq:monien}
\end{equation}
for $\det(\zeta(i+j+r))$. Monien himself notes that the passage from the discrete optimizer to a
continuous density is assumed there, and suggests Poissonization as the route to a fully rigorous
treatment; the case $r=0$ was subsequently justified through an explicit Riemann--Hilbert
construction. So the presence of an $N^{2}\log N$ scale in discrete Hankel determinants is not new,

Haynes and Zudilin \cite{HaynesZudilin2015} studied the arithmetic-progression matrices
$\det(\zeta(a(i+j)+b))$, which are the arithmetically interesting ones because for $a$ even and $b$
odd they avoid the known transcendental values $\zeta(2k)$. They state the expected law
\begin{equation}
\log\det_{1\le i,j\le n}\bigl(\zeta(a(i+j)+b)\bigr)=-a\,n^{2}\log n+O(n^{2}),
\label{eq:HZexpected}
\end{equation}
and then prove, by a short and completely elementary argument, the weaker one-sided estimate
\begin{equation}
\log\Bigl|\det_{1\le i,j\le n}\bigl(\zeta(a(i+j)+b)\bigr)\Bigr|\le-\frac a2\,n^{2}\log n+O(n^{2}).
\label{eq:HZproved}
\end{equation}
The distance between \eqref{eq:HZexpected} and \eqref{eq:HZproved} is a clean factor of two. For
the purpose of \cite{HaynesZudilin2015} --- proving one arithmetic statement by the shortest
available route --- nothing more was needed.

\label{ss:whytwo}
Their mechanism is integrality: if the $\zeta(an+b)$ are rational with common denominators $q_{n}$,
then $q_{n+1}\cdots q_{2n}H_{n}$ is a nonzero integer, hence at least $1$, so a \emph{small}
determinant forces \emph{large} denominators. Only an upper bound on $|H_{n}|$ can enter such an
argument. The conclusion is also coarser than the constant: what carries it is the divergence of
$\log n$, so $-\tfrac a2n^{2}\log n$ and $-a\,n^{2}\log n$ give the identical theorem
$q_{n}^{1/n}\to\infty$. Discarding the Vandermonde factor is then the natural move, and it is what
makes their proof fit in half a page.

This paper is downstream of the same representation --- $H_{n}$ written through row-linearity as a
weighted sum of Vandermonde products over reciprocal integers --- but does not proceed by sharpening
its estimates. It changes the question, and the change forces two-sided control: an estimate $\le-\kappa N^{2}\log N$
is compatible with every constant below $-\kappa$, so identification needs a lower bound as well,
and lower bounds are available here because the measures are positive, every term of the
Cauchy--Binet expansion is nonnegative, and one well-chosen term already suffices. That is the only
substantive ingredient used here and not there.

\subsection{The endpoint exponent}

Write the determinant as a Gram determinant of a positive discrete measure
$\mu=\sum_{j}w_{j}\delta_{x_{j}}$ whose atoms $x_{j}$ increase to a single accumulation endpoint,
normalized to $1$. Parametrize the atoms as $x_{j}=x(t_{j})$ along an increasing parameter
$t_{1}<t_{2}<\cdots\to\infty$, and define the \emph{endpoint exponent} $\beta>0$ by
\begin{equation}
1-x(t)=t^{-\beta}L(t),\qquad L \text{ slowly varying, i.e. } \log L(t)=o(\log t).
\label{eq:endpoint_RV}
\end{equation}
This is the natural hypothesis rather than an ad hoc one: it is the standard regular-variation
normal form for a monotone approach to an endpoint, and it is stable under exactly the
reparametrizations that occur in practice. In particular, if a counting law replaces $t_{j}\asymp j$
by $t_{j}\asymp j/\log j$, the induced factor $(\log j)^{2}$ is absorbed into $L$ and $\beta$ is
untouched.

\begin{theorem}[Endpoint-exponent law; informal statement]
Under \eqref{eq:endpoint_RV} and two coarse combinatorial conditions --- the node sequence is not
superlinearly dense on the logarithmic scale, and there exists a subfamily of $N$ atoms with
polynomially controlled height and separation --- the monomial Hankel determinant of $\mu$ satisfies
\[
\lim_{N\to\infty}\frac{\log D_{N}[\mu]}{N^{2}\log N}=-\beta .
\]
The same limit holds for the Gram determinant in any polynomial basis whose leading coefficients
$c_{n}$ satisfy $\sum_{n<N}\log|c_{n}|=O(N^{2})$.
\end{theorem}

The precise statement and proof are Theorem \ref{thm:endpoint} below. Three features are worth
isolating before the technical version.

\emph{It is a law, not an estimate.} The map $\beta\mapsto-\beta$ is a bijection between a geometric
invariant of the support and the leading determinant constant. Every particular evaluation
--- $\beta=1$ for $x_{m}=1/m$, $\beta=a$ for $x_{m}=m^{-a}$, $\beta=2$ for the Li--Cayley map --- is one
line. The triviality of an evaluation says nothing about the theorem, in the same way that
$\int x\dd x=x^{2}/2$ says nothing about the fundamental theorem of calculus.

\emph{It is two-sided and elementary.} The upper bound is the extremal property of monic orthogonal
polynomials against the competitor vanishing at the first $n$ atoms. The lower bound is a single
term of the Cauchy--Binet expansion, which is legitimate precisely because $\mu$ is positive, so all
terms of that expansion are nonnegative. No Riemann--Hilbert analysis, no equilibrium measure, no
lattice.

\emph{It says what cannot be in the constant.} Weights, provided they are subleading at scale
$N^{2}\log N$; the fine counting law; any arithmetic input --- none of these survive to leading order.
This is a negative statement of some force: if one hopes to see arithmetic in such a determinant,
one must look at order $N^{2}$ or below.

\subsection{What follows}

Section \ref{sec:zeta} applies the theorem to $\det(\zeta(a(i+j)+b))$. The nodes are $m^{-a}$, so
$\beta=a$, and all hypotheses hold with the sharpest possible error terms, giving
\eqref{eq:HZexpected} unconditionally and in the two-sided $O(n^{2})$ form. Corollary
\ref{cor:diophantine} converts this into the Diophantine statement of \cite{HaynesZudilin2015}
with twice the constant. Section \ref{sec:cayley} evaluates $\beta$ for the Li--Cayley map and
records the consequence $-2$. Section \ref{sec:numerics} reports high-precision computations, and
documents why local exponent estimators are the wrong figure of merit here. Section
\ref{sec:position} states plainly what in this paper was already known.

\section{The endpoint-exponent law}
\label{sec:law}

Let
\begin{equation}
\mu=\sum_{j\ge1}w_{j}\,\delta_{x_{j}},\qquad 0<w_{j},\qquad x_{j}=x(t_{j})\uparrow1,
\label{eq:measure}
\end{equation}
be a finite positive measure with distinct parameters $t_{1}<t_{2}<\cdots\to\infty$ and $x$
increasing and differentiable. Write
\begin{equation}
D_{N}[\mu]:=\det\Bigl(\int x^{r+s}\dd\mu(x)\Bigr)_{r,s=0}^{N-1}.
\label{eq:hankel}
\end{equation}

\begin{theorem}[Endpoint-exponent law]
\label{thm:endpoint}
Let $\beta>0$ and assume:
\begin{enumerate}[label=\textup{(\roman*)}]
\item \textup{(Endpoint power law.)} There are $c,C,t_{0}>0$ with, for $t\ge t_{0}$,
\begin{equation}
c\,t^{-\beta}\le1-x(t)\le C\,t^{-\beta},\qquad c\,t^{-\beta-1}\le x'(t)\le C\,t^{-\beta-1}.
\label{eq:power_law}
\end{equation}
\item \textup{(No superlinear logarithmic density.)}
\begin{equation}
\log t_{j}\ge(1-o(1))\log j .
\label{eq:node_scale}
\end{equation}
\item \textup{(A separated configuration.)} For every large $N$ there are support parameters
$s_{N,1}<\cdots<s_{N,N}$ and $\eta_{N}\to0$ with
\begin{equation}
s_{N,N}\le N^{1+\eta_{N}},\qquad s_{N,b}-s_{N,a}\ge(b-a)N^{-\eta_{N}}\quad(1\le a<b\le N).
\label{eq:separated}
\end{equation}
\item \textup{(Subleading weights.)}
\begin{equation}
\sum_{a=1}^{N}\log w(s_{N,a})\ge-o(N^{2}\log N).
\label{eq:weights}
\end{equation}
\end{enumerate}
Then
\begin{equation}
\boxed{\ \lim_{N\to\infty}\frac{\log D_{N}[\mu]}{N^{2}\log N}=-\beta.\ }
\label{eq:limit}
\end{equation}
The same limit holds for the Gram determinant in any polynomial basis $p_{0},p_{1},\dots$ with
$\deg p_{n}=n$ and leading coefficients $c_{n}$ obeying
\begin{equation}
\sum_{n=0}^{N-1}\log|c_{n}|=O(N^{2}).
\label{eq:basis}
\end{equation}
\end{theorem}

\begin{proof}
\emph{Upper bound.} Since $\mu$ is a finite positive measure with infinite support, the monomials
$1,x,\dots,x^{n}$ are linearly independent in $L^{2}(\mu)$: a nonzero polynomial has finitely many
zeros while the support has infinitely many points. Hence the monic orthogonal polynomial $\pi_{n}$
exists and is unique, and
\[
h_{n}:=\int|\pi_{n}|^{2}\dd\mu>0
\]
is the minimum of $\int|q|^{2}\dd\mu$ over monic $q$ of degree $n$. Triangular change of basis from
monomials to monic orthogonal polynomials gives the exact identity
\begin{equation}
D_{N}[\mu]=\prod_{n=0}^{N-1}h_{n}.
\label{eq:D_product}
\end{equation}

Take the monic competitor $q_{n}(x)=\prod_{j\le n}(x-x_{j})$. Since $x_{j}\uparrow1$, for $j>n$ and
$1\le r\le n$ we have $0<x_{j}-x_{r}\le1-x_{r}$, hence
$|q_{n}(x_{j})|\le\prod_{r\le n}(1-x_{r})$ for $j>n$. With $M:=\mu(\mathbb R)<\infty$,
\begin{equation}
h_{n}\le\int|q_{n}|^{2}\dd\mu=\sum_{j>n}w_{j}|q_{n}(x_{j})|^{2}
\le M\prod_{r=1}^{n}(1-x_{r})^{2}.
\label{eq:monic_upper}
\end{equation}

By (i), $\log(1-x_{j})\le\log C-\beta\log t_{j}$ for $j$ beyond a fixed index, and by (ii) there is
$\varepsilon_{j}\to0$ with $\log t_{j}\ge(1-\varepsilon_{j})\log j$, so
\begin{equation}
\log(1-x_{j})\le-\beta\log j+\beta\varepsilon_{j}\log j+O(1).
\end{equation}
The error sums to $o(n\log n)$: given $\epsilon>0$ choose $J$ with $|\varepsilon_{j}|\le\epsilon$ for
$j\ge J$, so $\sum_{J\le j\le n}|\varepsilon_{j}|\log j=O(\epsilon\, n\log n)$, and the finitely many
$j<J$ contribute $O_{J}(1)$. Since $\sum_{j\le n}\log j=n\log n-n+O(\log n)$,
\[
\sum_{j\le n}\log(1-x_{j})\le-\beta n\log n+o(n\log n),
\]
whence by \eqref{eq:monic_upper}
\begin{equation}
\log h_{n}\le-2\beta n\log n+o(n\log n).
\label{eq:h_upper}
\end{equation}
Summing over $0\le n<N$ preserves the scale: if $r_{n}=o(n\log n)$ then for every $\epsilon>0$ one
has $|r_{n}|\le\epsilon n\log n$ for large $n$, so $\sum_{n<N}r_{n}=o(N^{2}\log N)$. Using
$\sum_{n<N}n\log n=\tfrac12N^{2}\log N-\tfrac14N^{2}+O(N\log N)$ and \eqref{eq:D_product},
\begin{equation}
\log D_{N}[\mu]\le-\beta N^{2}\log N+o(N^{2}\log N).
\label{eq:det_upper}
\end{equation}

\emph{Lower bound.} Fix $N$ large, take the parameters $s_{N,1}<\cdots<s_{N,N}$ from (iii) and put
$y_{a}:=x(s_{N,a})$. The moment matrix generated by these selected atoms is $V_{N}W_{N}V_{N}^{T}$
with $(V_{N})_{r,a}=y_{a}^{r}$ and $W_{N}=\operatorname{diag}(w(s_{N,1}),\dots,w(s_{N,N}))$. The full
moment matrix is the sum of this positive semidefinite matrix and the contribution of the remaining
atoms; equivalently, in the Cauchy--Binet expansion of the full determinant every term is
nonnegative and the selected term occurs, so
\begin{equation}
D_{N}[\mu]\ \ge\ \Bigl(\prod_{a=1}^{N}w(s_{N,a})\Bigr)\prod_{1\le a<b\le N}(y_{b}-y_{a})^{2}.
\label{eq:one_term}
\end{equation}
This is the only place where positivity is used, and it is what makes a lower bound available at
all: for a signed measure the largest term of the expansion need not control the sum.

For $a<b$, monotonicity and (i) give
\[
y_{b}-y_{a}=\int_{s_{N,a}}^{s_{N,b}}x'(u)\dd u\ \ge\ c\,s_{N,N}^{-\beta-1}(s_{N,b}-s_{N,a}),
\]
since $u\le s_{N,N}$ on the interval and $u\mapsto u^{-\beta-1}$ decreases. By (iii),
\begin{equation}
y_{b}-y_{a}\ \ge\ c\,(b-a)\,N^{-\beta-1-(\beta+2)\eta_{N}}.
\label{eq:gap_lower}
\end{equation}
Taking logarithms and summing over the $\binom N2$ pairs,
\begin{equation}
2\sum_{a<b}\log(y_{b}-y_{a})\ \ge\ 2\sum_{a<b}\log(b-a)
-2\bigl[\beta+1+(\beta+2)\eta_{N}\bigr]\binom N2\log N+O(N^{2}).
\label{eq:vdm_lower}
\end{equation}
The identity $\prod_{a<b}(b-a)=\prod_{k<N}k!$ and Stirling give
\begin{equation}
\sum_{a<b}\log(b-a)=\tfrac12N^{2}\log N-\tfrac34N^{2}+O(N\log N),
\label{eq:pairsum}
\end{equation}
so, since $\eta_{N}\to0$,
\[
2\sum_{a<b}\log(y_{b}-y_{a})\ \ge\ -\beta N^{2}\log N-o(N^{2}\log N).
\]
Hypothesis (iv) says exactly that the selected weights do not degrade this at leading scale.
Inserting into \eqref{eq:one_term},
\begin{equation}
\log D_{N}[\mu]\ \ge\ -\beta N^{2}\log N-o(N^{2}\log N),
\label{eq:det_lower}
\end{equation}
which together with \eqref{eq:det_upper} proves \eqref{eq:limit}.

\emph{Polynomial bases.} Let $p_{n}(x)=c_{n}x^{n}+\cdots$ and let $C_{N}$ be the lower triangular
matrix taking $(1,x,\dots,x^{N-1})^{T}$ to $(p_{0},\dots,p_{N-1})^{T}$, so
$\det C_{N}=\prod_{n<N}c_{n}$. Then $G_{N}[p]=C_{N}H_{N}C_{N}^{T}$ exactly, hence
$\log\det G_{N}[p]=\log D_{N}[\mu]+2\sum_{n<N}\log|c_{n}|$, and \eqref{eq:basis} makes the second
term $O(N^{2})=o(N^{2}\log N)$.
\end{proof}

\subsection{The regularly varying form}

Hypothesis (i) is stated with two-sided constants for readability, but the proof uses it only
through $\log(1-x(t))=-\beta\log t+o(\log t)$ and through the integral lower bound for $x'$. Hence:

\begin{theorem}[Regular-variation form]
\label{thm:RV}
Suppose $1-x(t)=t^{-\beta}L(t)$ and $x'(t)=\beta\,t^{-\beta-1}L(t)\bigl(1+o(1)\bigr)$ with $L$ slowly
varying, i.e. $\log L(t)=o(\log t)$, and suppose \textup{(ii)--(iv)} of Theorem \ref{thm:endpoint}
hold. Then \eqref{eq:limit} holds.
\end{theorem}

\begin{proof}
In the upper bound replace $\log(1-x_{j})\le\log C-\beta\log t_{j}$ by
$\log(1-x_{j})=-\beta\log t_{j}+\log L(t_{j})$ and note $\sum_{j\le n}\log L(t_{j})=o(n\log n)$,
since $\log L(t)=o(\log t)$ and $\log t_{j}=O(\log j)$ on the relevant range. In the lower bound,
\eqref{eq:gap_lower} acquires the extra factor $L(s_{N,N})$, and
$\log L(N^{1+\eta_{N}})=o(\log N)$ contributes $o(N^{2}\log N)$ to \eqref{eq:vdm_lower}.
\end{proof}

\begin{remark}[Why regular variation is the right hypothesis, not an ad hoc one]
\label{rem:RV}
The condition $\log L(t)=o(\log t)$ is exactly what is needed for the exponent to be well defined,
and it is exactly what makes the statement invariant under the reparametrizations that occur when a
node sequence is indexed by height rather than by rank. If the nodes are indexed by a counting law
with $t_{j}\asymp j/\log j$, then
$1-x(t_{j})\asymp j^{-\beta}(\log j)^{\beta}$: the induced factor $(\log j)^{\beta}$ is slowly
varying, so it is absorbed by $L$ and leaves $\beta$ invariant. Any hypothesis stated with
$1-x(t_{j})\asymp j^{-\beta}$ and hard constants would fail here, and would have to be repaired case
by case. This is the technical content of the phrase ``the counting law cannot appear in the leading
coefficient.''
\end{remark}

\begin{remark}[What is universal]
\label{rem:universal}
Theorem \ref{thm:endpoint} shows the leading constant is the endpoint exponent and nothing else.
The node counting law enters only through the verification that enough atoms exist at the right
logarithmic scale; $t_{j}\asymp j$, $t_{j}\asymp j/\log j$, and all compatible logarithmic
refinements give the same $-\beta$. Weights enter only through the one-sided condition (iv).
Consequently arithmetic information, if it is present at all, can first appear at order $N^{2}$.
\end{remark}

\section{Hankel determinants of zeta values: the expected law, unconditionally}
\label{sec:zeta}

Fix positive integers $a$ and $b\ge0$ and set
$H_{n}^{(a,b)}:=\det_{1\le i,j\le n}\bigl(\zeta(a(i+j)+b)\bigr)$. Expanding each entry,
\begin{equation}
\zeta(a(i+j)+b)=\sum_{m\ge1}m^{-b}\bigl(m^{-a}\bigr)^{i+j},
\end{equation}
so $H_{n}^{(a,b)}=D_{n}[\nu]$ is the monomial Hankel determinant \eqref{eq:hankel} of the positive
discrete measure
\begin{equation}
\nu=\sum_{m\ge1}m^{-(b+2a)}\,\delta_{u_{m}},\qquad u_{m}=m^{-a},
\label{eq:nu}
\end{equation}
in indices $r,s=0,\dots,n-1$. Two immediate consequences of positivity: $H_{n}^{(a,b)}>0$ for every
$n$, so no nonvanishing argument is required; and the Cauchy--Binet lower bound of Theorem
\ref{thm:endpoint} is available.

The atoms of $\nu$ accumulate at $0$ rather than at $1$. Apply the reflection $x\mapsto1-x$: the
change of basis from $\{1,x,\dots,x^{N-1}\}$ to $\{1,(1-x),\dots,(1-x)^{N-1}\}$ is triangular with
diagonal entries $\pm1$, so the Gram determinants coincide. The reflected measure has atoms
$x_{m}=1-m^{-a}\uparrow1$ with support map
\begin{equation}
x(t)=1-t^{-a},\qquad 1-x(t)=t^{-a},\qquad x'(t)=a\,t^{-a-1}.
\end{equation}

\begin{theorem}[Full leading asymptotics of $H_{n}^{(a,b)}$]
\label{thm:zeta}
For all positive integers $a$ and all integers $b\ge0$,
\begin{equation}
\boxed{\ \log\det_{1\le i,j\le n}\bigl(\zeta(a(i+j)+b)\bigr)=-a\,n^{2}\log n+O(n^{2}).\ }
\label{eq:zeta_law}
\end{equation}
In particular $\log H_{n}^{(a,b)}/(n^{2}\log n)\to-a$. The statement is unconditional.
\end{theorem}

\begin{proof}
Hypothesis (i) of Theorem \ref{thm:endpoint} holds with $\beta=a$ and constants $c=C=1$ for the
first pair and $c=C=a$ for the second. Hypothesis (ii) holds with $t_{j}=j$ and no error at all.
Hypothesis (iii) holds with $s_{N,a}=a$ and $\eta_{N}=0$: then $s_{N,N}=N$ and
$s_{N,b}-s_{N,a}=b-a$. Hypothesis (iv) holds because
$\sum_{m\le N}\log m^{-(b+2a)}=-(b+2a)\,N\log N+O(N)=o(N^{2}\log N)$.

For the sharper error term, run the two bounds of Theorem \ref{thm:endpoint} with these exact data.
In the upper bound, $\log(1-x_{j})=-a\log j$ exactly, so
$\sum_{j\le n}\log(1-x_{j})=-a\,n\log n+O(n)$ and $\log h_{n}\le-2a\,n\log n+O(n)$; summing gives
$\log H_{n}^{(a,b)}\le-a\,n^{2}\log n+O(n^{2})$. In the lower bound, $\eta_{N}=0$ turns
\eqref{eq:gap_lower} into $y_{b}-y_{a}\ge a(b-a)N^{-a-1}$, and \eqref{eq:vdm_lower} together with
\eqref{eq:pairsum} gives $-a\,n^{2}\log n+O(n^{2})$; the weight term is $O(n\log n)$.
\end{proof}

\begin{remark}[Relation to \cite{Monien2009} and \cite{HaynesZudilin2015}]
\label{rem:position}
For $a=1$, \eqref{eq:zeta_law} reproduces the leading term of \eqref{eq:monien} and is consistent
with it; Monien's calculation is finer at order $n^{2}$ (it produces the explicit constant
$\log 2-\tfrac32$) but is carried out for a single model and, in the general case, at a variational
level the author himself flags as not fully rigorous. Theorem \ref{thm:zeta} is coarser at order
$n^{2}$ and broader: it holds for every arithmetic progression, unconditionally, and it is two-sided.

Against \cite{HaynesZudilin2015} the comparison is exact and there is no contradiction. They prove
\eqref{eq:HZproved}, an upper bound with coefficient $-a/2$, and record $-a$ as expected; Theorem
\ref{thm:zeta} supplies $-a$ in both directions, and their Theorem 1 is unaffected. In the present
framework the factor two is located precisely: the leading behaviour is a product of two
contributions of equal size, a weight contribution and a Vandermonde contribution, each
$\exp(-\tfrac a2n^{2}\log n)$, and the estimate \eqref{eq:HZproved} retains the first and discards
the second. Retaining the second requires positivity, which has no role in a one-sided bound; see
\S\ref{ss:whytwo}.
\end{remark}

\begin{corollary}[Diophantine consequence, with twice the constant]
\label{cor:diophantine}
Let $a,b$ be positive integers and suppose all the numbers $\zeta(an+b)$, $n\ge1$, are rational.
Let $q_{n}$ be the common denominator of $\zeta(a+b),\dots,\zeta(an+b)$. Then
\begin{equation}
\liminf_{n\to\infty}\frac{\log q_{n}}{n\log n}\ \ge\ \frac a2 .
\end{equation}
Consequently, either $\zeta(an+b)$ is irrational for infinitely many $n$, or $q_{n}$ grows at least
like $n^{(a/2+o(1))n}$. The same argument fed with \eqref{eq:HZproved} yields only $a/4$.
\end{corollary}

\begin{proof}
As in \cite{HaynesZudilin2015}, $q_{k}\mid q_{n}$ for $k\le n$ and
\[
q_{n+1}q_{n+2}\cdots q_{2n}\,H_{n}^{(a,b)}=\det_{1\le i,j\le n}\bigl(q_{i+n}\zeta(a(i+j)+b)\bigr)\in\mathbb Z .
\]
By positivity $H_{n}^{(a,b)}>0$ for every $n$, so the integer is nonzero and
$q_{n+1}\cdots q_{2n}\,H_{n}^{(a,b)}\ge1$ for all $n$. Hence
$\log(q_{n+1}\cdots q_{2n})\ge-\log H_{n}^{(a,b)}=a\,n^{2}\log n+O(n^{2})$ by Theorem \ref{thm:zeta}.
Since $q_{k}\le q_{2n}$ for $k\le2n$, the left side is at most $n\log q_{2n}$, so
$\log q_{2n}\ge a\,n\log n+O(n)$. Writing $N=2n$ gives
$\log q_{N}\ge\tfrac a2N\log N+O(N)$. Replacing \eqref{eq:zeta_law} by \eqref{eq:HZproved}
halves the right-hand side throughout.
\end{proof}

\begin{remark}
Corollary \ref{cor:diophantine} does not decide the irrationality of any odd zeta value, and no such
claim is made, and the qualitative conclusion $q_{n}^{1/n}\to\infty$ of \cite{HaynesZudilin2015}
follows equally from either constant. It is recorded for one narrow purpose: it is the only place
in this paper where the difference between $-a/2$ and $-a$ produces a different sentence, and it
therefore certifies that the sharp constant is a real quantity rather than a cosmetic one.
\end{remark}

\section{The Li--Cayley exponent: $\beta=2$}
\label{sec:cayley}

The Cayley-type map underlying Li's criterion \cite{Li1997,BombieriLagarias1999} sends a nontrivial
zero $\rho=\tfrac12+i\gamma$ of $\zeta$ to $z_{\rho}=1-1/\rho$, which lies on the unit circle exactly
when $\rho$ is on the critical line. Its real part
\begin{equation}
x_{\gamma}:=\operatorname{Re}\Bigl(1-\frac1\rho\Bigr)=\frac{4\gamma^{2}-1}{4\gamma^{2}+1}\in(-1,1)
\label{eq:cayley_map}
\end{equation}
is the natural real coordinate, and
\begin{equation}
1-x_{\gamma}=\frac{2}{4\gamma^{2}+1},
\qquad\text{so}\qquad
\log(1-x_{\gamma})=-2\log\gamma+O(\gamma^{-2}).
\label{eq:cayley_endpoint}
\end{equation}

\begin{proposition}
\label{prop:beta2}
The map \eqref{eq:cayley_map}, viewed as $x(t)=(4t^{2}-1)/(4t^{2}+1)$, satisfies
$1-x(t)\asymp t^{-2}$ and $x'(t)\asymp t^{-3}$. Its endpoint exponent is $\beta=2$.
\end{proposition}

\begin{proof}
Immediate from \eqref{eq:cayley_endpoint} and $x'(t)=16t/(4t^{2}+1)^{2}$.
\end{proof}

Note the shape of \eqref{eq:cayley_endpoint}: the exponent is exactly $2$, with no logarithmic
correction, and its derivation does not invoke the Riemann--von Mangoldt formula at all. That
formula enters only when one insists on indexing by rank rather than by height, and Remark
\ref{rem:RV} explains why that reindexing cannot move $\beta$.

The two combinatorial hypotheses are supplied unconditionally by standard zero-counting.

\begin{lemma}[Coarse location of distinct ordinates]
\label{lem:coarse}
Let $\gamma_{1}^{*}<\gamma_{2}^{*}<\cdots$ be the distinct positive ordinates. There are $c>0$ and
$j_{0}$ with $\gamma_{j}^{*}\ge c\,j/\log j$ for $j\ge j_{0}$. No simplicity hypothesis is needed.
\end{lemma}

\begin{proof}
Let $N^{*}(T)$ count distinct positive ordinates and $N(T)$ count them with multiplicity, so
$N^{*}(T)\le N(T)$. Riemann--von Mangoldt gives $N(T)\ll T\log T$, hence
$j\le N(\gamma_{j}^{*})\ll\gamma_{j}^{*}\log\gamma_{j}^{*}$. Multiple zeros only decrease $N^{*}$
and so move the distinct ordinates upward; the bound is robust in multiplicity.
\end{proof}

\begin{lemma}[A deterministically separated subsequence]
\label{lem:separated}
There are $H>0$, $T_{0}>0$ and distinct positive ordinates $\widetilde\gamma_{j}$ with
$T_{0}+2Hj\le\widetilde\gamma_{j}\le T_{0}+2Hj+H$ for all $j\ge1$. Hence
$\widetilde\gamma_{j}\asymp j$ and $\widetilde\gamma_{b}-\widetilde\gamma_{a}\ge H(b-a)$.
\end{lemma}

\begin{proof}
Write $N(T)=F(T)+R(T)$ with $F(T)=\frac T{2\pi}\log\frac T{2\pi}-\frac T{2\pi}$ and
$R(T)=O(\log T)$, say $|R(T)|\le A\log T$ for $T\ge T_{1}$. For fixed $H\ge1$,
$F(T+H)-F(T)=\frac H{2\pi}\log T+O_{H}(1)$. Choosing $H>4\pi A$ gives
$N(T+H)-N(T)\ge(\frac H{2\pi}-2A+o(1))\log T>0$ for large $T$, so every interval $[T,T+H]$ with
$T\ge T_{0}$ contains an ordinate. Apply this to the disjoint boxes $[T_{0}+2Hj,T_{0}+2Hj+H]$.
\end{proof}

\begin{corollary}[Universal Cayley constant]
\label{cor:cayley}
Let $\mu$ be any positive discrete measure supported on the Cayley images \eqref{eq:cayley_map} of
the ordinates of the nontrivial zeros of $\zeta$, with weights satisfying \eqref{eq:weights} along
the separated subsequence of Lemma \ref{lem:separated}, and let $D_{N}$ be its Gram determinant in
any polynomial basis obeying \eqref{eq:basis}. Then
\begin{equation}
\boxed{\ \lim_{N\to\infty}\frac{\log D_{N}}{N^{2}\log N}=-2.\ }
\end{equation}
Moreover the value $-2$ is unchanged if the $T\log T$ counting law is replaced by a $T$ counting
law, or by any logarithmic refinement compatible with \eqref{eq:node_scale} and
\eqref{eq:separated}.
\end{corollary}

\begin{proof}
Proposition \ref{prop:beta2} gives $\beta=2$; Lemma \ref{lem:coarse} gives \eqref{eq:node_scale};
Lemma \ref{lem:separated} gives \eqref{eq:separated} with $\eta_{N}=0$ up to constants. Apply
Theorem \ref{thm:endpoint}. The invariance statement is Remark \ref{rem:universal}.
\end{proof}

\begin{remark}[Where the Riemann--von Mangoldt refinement can first appear]
Quantitatively, the upper bound in this setting carries a correction: from Lemma \ref{lem:coarse},
$1-x_{\gamma_{j}^{*}}\ll(\log j)^{2}/j^{2}$, which propagates to
\[
\log D_{N}\le-2N^{2}\log N+2N^{2}\log\log N+O(N^{2}),
\]
while the Cauchy--Binet lower bound gives $-2N^{2}\log N+O(N^{2})$. The $N^{2}\log\log N$ term is not
declared an artifact here: positivity guarantees that one favourable subset bounds the determinant
from below, but not that the logarithm of the whole sum equals the logarithm of its largest term, so
an entropy contribution may survive below the leading scale. Both bounds have the same leading
coefficient $-2$, which is all that Corollary \ref{cor:cayley} asserts.
\end{remark}

\begin{remark}[Automorphic $L$-functions]
Lagarias \cite{Lagarias2007} defines Li coefficients for principal automorphic $L$-functions and
relates them to Weil's quadratic functional. Whenever a member of that class, under GRH, produces a
positive Cayley zero measure with the same coarse counting and separation properties, Corollary
\ref{cor:cayley} returns the same constant $-2$: degree and conductor may alter lower-order terms
but not the endpoint exponent. This is a conditional universality statement, and the present paper
does not construct the relevant measure for every such $L$-function.
\end{remark}

\section{Numerics, and a warning about local exponents}
\label{sec:numerics}

\subsection{Confirmation of $-a$}

Table \ref{tab:zeta} reports $\log H_{n}^{(a,b)}$ computed at $2500$ decimal digits for
$6\le n\le32$, together with the coefficient $c$ obtained by least squares in the model
\begin{equation}
\log H_{n}^{(a,b)}=c\,n^{2}\log n+d\,n^{2}+e\,n\log n+f\,n+g .
\label{eq:fit_model}
\end{equation}

\begin{table}[htbp]
\centering
\begin{tabular}{ccrrrr}
\toprule
$a$ & $b$ & $\dfrac{\log H_{32}}{32^{2}\log32}$ & $\dfrac{\log H_{32}+a\cdot32^{2}\log32}{32^{2}}$ & fitted $c$ & fitted $d$\\
\midrule
1 & 0 & $-0.7961$ & $0.7067$ & $-0.9990$ & $0.8010$\\
1 & 1 & $-0.8248$ & $0.6071$ & $-0.9984$ & $0.7975$\\
2 & 1 & $-1.3962$ & $2.0925$ & $-1.9949$ & $2.4870$\\
3 & 2 & $-2.0191$ & $3.3996$ & $-2.9905$ & $4.0725$\\
\bottomrule
\end{tabular}
\caption{Hankel determinants of zeta values. Column 3 is the raw ratio, which is
still far from $-a$ at $n=32$ because of the $O(n^{2})$ term; column 4 shows that the residual
after subtracting $-a\,n^{2}\log n$ is indeed of size $n^{2}$; column 5 is the fitted leading
coefficient from \eqref{eq:fit_model}, agreeing with $-a$ to three digits. All values computed at
$2500$ digits.}
\label{tab:zeta}
\end{table}

Two points deserve emphasis. The raw ratio $\log H_{32}/(32^{2}\log32)$ equals $-0.796$ when the
true constant is $-1$: at $n=32$ the naive figure of merit is off by $20\%$, purely because
$\log32\approx3.47$ is not large. And the residual column is flat in $n$, confirming the $O(n^{2})$
error term of Theorem \ref{thm:zeta} rather than merely $o(n^{2}\log n)$.

\subsection{The weights really do not matter}

Theorem \ref{thm:endpoint} asserts that neither the size of the weights, subject to
\eqref{eq:weights}, nor the removal of most of the atoms can move the constant. Both are testable.
Taking $w_{m}=m^{-b}$ supported on $m\equiv1\pmod4$ replaces $\zeta$ by $\zeta+L(\cdot,\chi_{4})$ in
the matrix entries and deletes three quarters of the support; taking $b=9$ instead of $b=1$ makes
the weights decay far faster. Fitting \eqref{eq:fit_model} on $6\le n\le20$ at $700$ digits:

\begin{center}
\begin{tabular}{llr}
\toprule
$a$ & measure & fitted $c$\\
\midrule
$1$ & $w_{m}=m^{-1}$, all atoms & $-0.997$\\
$1$ & $w_{m}=m^{-1}$, atoms $m\equiv1\ (4)$ & $-0.931$\\
$2$ & $w_{m}=m^{-1}$, atoms $m\equiv1\ (4)$ & $-2.328$\\
$1$ & $w_{m}=m^{-9}$, all atoms & $-0.961$\\
\bottomrule
\end{tabular}
\end{center}

All four are consistent with $-a$ at the resolution available from $n\le20$; the calibration row
shows that the plain case is already at $-0.997$ there, and the deviations in the sparse rows are
what one expects, since deleting atoms shifts $\log(1-x_{k})$ by the constant $-a\log4$ and
therefore inflates the $n^{2}$ coefficient, which a fit truncated at $n=20$ cannot cleanly separate
from $n^{2}\log n$. These are consistency checks, not sharp confirmations.

\subsection{Local exponents are the wrong figure of merit}

For the Li--Cayley nodes the situation is worse, and instructively so. Write
$S(M):=2\sum_{a\le M}\log\delta_{a}+2\sum_{a<b\le M}\log(\delta_{a}-\delta_{b})$ with
$\delta_{j}=2/(4\gamma_{j}^{2}+1)$, the logarithm of the Cauchy--Binet term that Theorem
\ref{thm:endpoint} shows to be leading. Using the first $140$ ordinates and, beyond that, the
Riemann--von Mangoldt/Lambert approximation to $\gamma_{j}$, one finds

\begin{center}
\begin{tabular}{lrrrrr}
\toprule
$M$ & $200$ & $600$ & $1400$ & $2600$ & $4000$\\
\midrule
$S(M)/(M^{2}\log M)$ & $-2.148$ & $-2.014$ & $-1.951$ & $-1.919$ & $-1.902$\\
\bottomrule
\end{tabular}
\end{center}

The local ratio does not settle on $-2$; it drifts. The reason is exactly the slowly varying factor
of Remark \ref{rem:RV}. With $\gamma_{j}\asymp j/\log j$ one has
$\log\delta_{j}=-2\log j+2\log\log j+O(1)$, so
\[
\frac{S(M)}{M^{2}\log M}=-2+\frac{2\log\log M}{\log M}+O\Bigl(\frac1{\log M}\Bigr),
\]
and the correction decays like $1/\log M$: at $M=4000$ it is still of size $0.5$. Fitting
$\log S$ against a basis that includes an $M^{2}\log\log M$ term restores the constant
(e.g.\ $-2.09$ for the synthetic model $\delta_{j}=j^{-2}(\log j)^{2}$ at $M\le600$, against
$-1.69$ for the naive basis; and $-2.00$ for $\delta_{j}=j^{-2}$, where no correction exists).

The moral is methodological and worth stating flatly: an empirical exponent extracted from a
finite range of $j$ measures the endpoint exponent \emph{plus} the local logarithmic derivative of
$L$. It is not an estimator of $\beta$ at any accessible range, and a numerically slow approach to
$-2$ is a property of the parametrization, not evidence against the constant. Theorem
\ref{thm:endpoint} is, in this regime, the only reliable access to the coefficient.

\section{What is new and what is not}
\label{sec:position}

\subsection{The law in Jacobi form, and what it actually asserts}
\label{ss:jacobi}

It is useful to know what Theorem \ref{thm:endpoint} is equivalent to, because the equivalent form
makes its content explicit and shows what has to be true for it to hold.

\begin{proposition}
\label{prop:jacobi}
Let $\mu$ be a probability measure as in \eqref{eq:measure}, let $p_{n}$ be its orthonormal
polynomials with three-term recurrence $xp_{n}=b_{n}p_{n+1}+a_{n}p_{n}+b_{n-1}p_{n-1}$, and let
$D_{N}[\mu]$ be as in \eqref{eq:hankel}. Then, identically,
\begin{equation}
\log D_{N}[\mu]=2\sum_{m=0}^{N-2}(N-1-m)\log b_{m}.
\label{eq:jacobi_identity}
\end{equation}
Consequently the conclusion \eqref{eq:limit} of Theorem \ref{thm:endpoint} is equivalent to
\begin{equation}
\frac{2}{N^{2}}\sum_{m<N}(N-m)\log b_{m}=-\beta\log N+o(\log N),
\label{eq:jacobi_equivalent}
\end{equation}
that is, to $b_{m}=m^{-\beta+o(1)}$ in the logarithmic Ces\`aro sense.
\end{proposition}

\begin{proof}
The monic orthogonal polynomial is $\pi_{n}=p_{n}/k_{n}$ where $k_{n}$ is the leading coefficient of
$p_{n}$, and the recurrence gives $k_{n}=k_{n-1}/b_{n-1}$ with $k_{0}=1$, so
$k_{n}=(b_{0}\cdots b_{n-1})^{-1}$ and $h_{n}=k_{n}^{-2}=(b_{0}\cdots b_{n-1})^{2}$. Insert this in
\eqref{eq:D_product} and collect powers. For the equivalence, substitute
$\log b_{m}=-\beta\log m+r_{m}$ into \eqref{eq:jacobi_identity} and use
$\sum_{m<N}(N-m)\log m=\tfrac12N^{2}\log N+O(N^{2})$.
\end{proof}

So the theorem says: \emph{the endpoint exponent of the support controls the decay exponent of the
Jacobi off-diagonal coefficients, on logarithmic average, uniformly over admissible weights.} This
is an implication from the geometry of the atoms to the recurrence coefficients. The moment-problem
literature closest to it --- \cite{BergPedersen1994} on the order and type of indeterminate Hamburger
problems, and Berg--Szwarc, who identify the order with the exponent of convergence of $(b_{n})$
\cite{BergSzwarc2014} --- runs predominantly in the opposite direction, taking regularity of
$(b_{n})$ as hypothesis. Theorem \ref{thm:endpoint} never touches the recurrence at all: it reaches
$D_{N}$ through an extremal competitor and a single Cauchy--Binet term, and \eqref{eq:jacobi_identity}
is then a translation.

\subsection{Where the difficulty lies, and where it does not}
\label{ss:trivial}

The argument of \S\ref{sec:law} is short, and every ingredient in it is classical: the factorization
$D_{N}=\prod_{n<N}h_{n}$, the extremal characterization of $h_{n}$, the positivity of the
Cauchy--Binet expansion. These tools have been available for decades; what had not been available
is the question. The brevity is moreover not incidental to the answer but a condition of it. An
argument requiring precise local information about the support would be unusable on the ordinates of
$\zeta$, about which such information is not available; only hypotheses as coarse as
\eqref{eq:node_scale} and \eqref{eq:separated} can be verified there, and coarse hypotheses admit
only short proofs.

The substance of this paper is therefore not a technical barrier that could not be passed. It is
that a regime the standard theory leaves undescribed turns out to carry a clean invariant, that the
invariant is computable from the transform alone, and that useful things follow from knowing it
--- \S\ref{sec:conclusions}.

\emph{That the standard theory leaves it undescribed is worth being precise about.} For a measure on
a compact set $E$ of positive logarithmic capacity, regular in the
Stahl--Totik sense, one has $h_{n}^{1/n}\to\operatorname{cap}(E)^{2}$, hence
$\log D_{N}\sim N^{2}\log\operatorname{cap}(E)$: the $N^{2}$ law with a capacity constant
\cite{StahlTotik1992,Totik2005Orthogonal}. The supports considered here are countable, so
$\operatorname{cap}(E)=0$, and that law degenerates: it asserts only $h_{n}^{1/n}\to0$, with no
rate. Theorem \ref{thm:endpoint} operates exactly in the gap this leaves, supplying the rate
$\log h_{n}\sim-2\beta n\log n$ and showing that it is governed by a single number. It is therefore
not a specialization of Szeg\H o--Widom or of Stahl--Totik regularity; those give nothing to
specialize.

\emph{Two points in it are not bookkeeping.} First, the lower bound retains one term of an infinite
nonnegative sum, and a priori the remaining terms could contribute an entropy at the same order
$N^{2}\log N$; that they do not is established by the matching upper bound, not assumed. Second, the
hypotheses are designed rather than inherited: the regularly varying form of Theorem \ref{thm:RV} is
what allows the zeta ordinates to be handled with no appeal to Riemann--von Mangoldt beyond the
coarse Lemmas \ref{lem:coarse} and \ref{lem:separated}, and a formulation with hard constants in the
index would not survive the reparametrization (Remark \ref{rem:RV}).

The summary is that the result is short once the question is posed, and that the question had not
been posed. Short and new are compatible; classical and new are not, and this is the former.

\subsection{Search for prior appearances}
\label{ss:priorart}

We searched for the statement of Theorem \ref{thm:endpoint}, or for any identification of the
leading Hankel constant with a support exponent, in the following literatures, and did not find it:
Hankel determinant asymptotics by Riemann--Hilbert methods, including Fisher--Hartwig singularities
and one-cut regular potentials; discrete orthogonal polynomial ensembles, where
\cite{Johansson2001} and its successors are confined to uniform lattices, the saturation constraint
$\rho\le1$ being an expression of unit spacing, whereas the supports here are neither uniform nor
lattice-like; discrete orthogonal polynomials on infinite node sets with an accumulation point,
where \cite{WuLinXuZhao2016} performs a full uniform analysis of one specific model; the theory of
indeterminate moment problems, where \cite{BergSzwarc2014} and the order-and-type literature give
the closest invariants, as discussed in \S\ref{ss:jacobi}; and general orthogonal polynomial theory
on sets of zero capacity \cite{StahlTotik1992,Totik2005Orthogonal}. We also examined the citation
records of \cite{Monien2009} and \cite{HaynesZudilin2015}: neither has been followed up on the
asymptotic question.

A negative search is not a proof of novelty, and we state the claim in that form. What we assert is
that the statement is not a corollary of any framework we have been able to locate, and that in the
one place where it overlaps a published estimate --- \eqref{eq:HZproved} --- it gives the constant
that estimate does not.

\subsection{Summary}

For the avoidance of doubt, the following are \emph{not} claimed as novelties.

\begin{itemize}
\item The $N^{2}\log N$ decay scale of discrete Hankel determinants built from Dirichlet series.
This is in \cite{Monien2009} and is recorded in \cite{HaynesZudilin2015}.
\item The discrete Vandermonde / Cauchy--Binet representation of such determinants, which is
\cite{Monien2009} and is also the engine of the elementary bound in \cite{HaynesZudilin2015}.
\item The identity $D_{N}[\mu]=\prod_{n<N}h_{n}$ and the extremal characterization of $h_{n}$, which
are textbook facts about orthogonal polynomials \cite{Chihara1978}.
\item Discrete-Coulomb-gas and constrained-equilibrium methods, which are Johansson's
\cite{Johansson2001}; nor the asymptotic theory of discrete orthogonal polynomials on infinite node
sets with an accumulation point \cite{WuLinXuZhao2016}, or the general theory of orthogonal
polynomials on sets of zero capacity \cite{StahlTotik1992,Totik2005Orthogonal}.
\item The Li--Cayley map itself \cite{Li1997,BombieriLagarias1999} and its automorphic extension
\cite{Lagarias2007}.
\end{itemize}

The following are the paper's contributions.

\begin{itemize}
\item \textbf{The law $\beta\mapsto-\beta$.} Theorems \ref{thm:endpoint} and \ref{thm:RV} identify
the leading determinant constant with a single geometric invariant of the support --- the endpoint
exponent --- under hypotheses that impose no lattice, no explicit node formula and no weight
asymptotics. The statement is two-sided and the proof is short; \S\ref{ss:trivial} locates the
contribution and \S\ref{ss:priorart} records the search for prior appearances.
\item \textbf{The zeta-value asymptotic in full.} Theorem \ref{thm:zeta} proves
\eqref{eq:HZexpected} unconditionally, in two-sided $O(n^{2})$ form, for every arithmetic
progression, without the Riemann--Hilbert apparatus needed for the single case $r=0$. This is the
law recorded as expected in \cite{HaynesZudilin2015}, and it serves here as the test of the general
principle on the one family where the right constant was already anticipated.
\item \textbf{A quantitative Diophantine by-product.} Corollary \ref{cor:diophantine} doubles the
denominator-growth constant obtainable by the method of \cite{HaynesZudilin2015}, without altering
the qualitative conclusion they draw.
\item \textbf{The evaluation $\beta_{\mathrm{Li\text-Cayley}}=2$.} Proposition \ref{prop:beta2} is a
one-line computation, and that is the point: the content is in the theorem, not the evaluation. Its
consequence, Corollary \ref{cor:cayley}, is a rigidity statement --- the leading determinant
behaviour of every positive Gram system on Cayley images of zeta ordinates is $-2$, and the
$T\log T$ counting law cannot reach it.
\item \textbf{A negative structural conclusion.} Because the constant is purely geometric, any
arithmetic signal in these determinants lives at order $N^{2}$ or below. Remark \ref{rem:universal}
makes this precise, and Section \ref{sec:numerics} shows why the leading order is nonetheless hard
to see numerically.
\end{itemize}

The discrete Vandermonde representation and the correct scale are due to \cite{Monien2009} and
\cite{HaynesZudilin2015}; everything below the leading order in \eqref{eq:monien} is beyond the
present methods and is not attempted. What is claimed is the change of question and its answer.

\section{Conclusions: what the law is for}
\label{sec:conclusions}

The endpoint-exponent law is a statement about a leading constant, and leading constants are cheap
to state and easy to dismiss. This section sets out, concretely, what follows from knowing that the
constant is $-\beta$ and nothing else. Two of the consequences are quantitative and immediately
usable; one is negative and, we think, the most valuable; one reframes a familiar phenomenon in
numerical linear algebra.

\subsection{A quantified precision barrier}
\label{ss:conditioning}

\begin{proposition}
\label{prop:kappa}
Let $\mu$ satisfy the hypotheses of Theorem \ref{thm:endpoint}, let $G_{N}$ be the $N\times N$ Gram
matrix of \eqref{eq:hankel}, and let $\kappa_{2}(G_{N})=\lambda_{\max}/\lambda_{\min}$ be its
spectral condition number. Then
\begin{equation}
\liminf_{N\to\infty}\frac{\log\kappa_{2}(G_{N})}{N\log N}\ \ge\ \beta .
\label{eq:kappa}
\end{equation}
\end{proposition}

\begin{proof}
$G_{N}$ is positive definite, so $\lambda_{\min}^{N}\le\prod_{i}\lambda_{i}=D_{N}[\mu]$, giving
$\log\lambda_{\min}\le N^{-1}\log D_{N}[\mu]=-\beta N\log N+o(N\log N)$ by \eqref{eq:limit}.
Also $\lambda_{\max}\ge(G_{N})_{00}=\mu(\mathbb R)=:M>0$. Hence
$\log\kappa_{2}\ge\log M-N^{-1}\log D_{N}[\mu]=\beta N\log N+o(N\log N)$.
\end{proof}

The practical content is a lower bound on the working precision required to compute anything
spectral about $G_{N}$ --- a Cholesky factorization, a smallest eigenvalue, a pivot sign. Roughly
$\beta N\log N/\log10$ significant decimal digits are needed merely to separate the extreme
eigenvalues. For the Li--Cayley value $\beta=2$:

\begin{center}
\begin{tabular}{lrrrrr}
\toprule
$N$ & $50$ & $100$ & $200$ & $500$ & $1000$\\
\midrule
decimal digits & $\approx170$ & $\approx400$ & $\approx920$ & $\approx2700$ & $\approx6000$\\
\bottomrule
\end{tabular}
\end{center}

This is not an empirical remark about any particular implementation. It is a proved lower bound
that no algorithm can evade, since it is a property of the matrix and not of the method. Any
programme that proposes to certify a positivity property of such a family numerically can read off
its own cost before writing any code, and can recognize in advance the size at which the
computation stops meaning anything in a given precision. In our experience this is the single most
frequently useful consequence of the theorem.

\subsection{A no-go: arithmetic cannot reach the leading order}
\label{ss:nogo}

By Theorem \ref{thm:endpoint} the constant $-\beta$ is a function of the support geometry alone.
Nothing else survives to order $N^{2}\log N$: not the zero-counting law, not the weights beyond the
one-sided condition \eqref{eq:weights}, not any lattice or spacing structure, and in particular no
arithmetic input whatsoever. Remark \ref{rem:universal} states this, and Corollary
\ref{cor:cayley} exhibits it: replacing the Riemann--von Mangoldt law $T\log T$ by a bare counting
law $T$ leaves the constant at $-2$.

Two consequences follow, and they point in the same direction.

First, a methodological warning. An observed agreement between a computed determinant and its
expected leading decay is \emph{not} evidence for any arithmetic hypothesis. It is a consequence of
the support being where it is, and it would persist under substantial deformations of the zero
configuration. A numerical check at leading order confirms a universality statement wearing
arithmetic clothing.

Second, a redirection of effort. If arithmetic information is present in these determinants at all,
it lives at order $N^{2}$ or below --- exactly the order that Monien's variational analysis reaches
for one model \eqref{eq:monien} and that our methods deliberately do not. Any programme hoping to
extract spectral or arithmetic content from the bulk asymptotics of Li--Weil type matrices must be
posed at that order. Section \ref{sec:numerics} shows that this is a demanding target: even the
leading term is hard to see numerically, and the subleading one sits underneath an
$N^{2}\log\log N$ correction whose status we have deliberately left open.

We regard this as the most useful output of the paper. It closes an avenue, with proof rather than
intuition, and it says where to look instead.

\subsection{$\beta$ as a design parameter}

Since $\beta$ is read off from the transform and $\beta$ alone controls both the determinant
constant \eqref{eq:limit} and the conditioning bound \eqref{eq:kappa}, the endpoint exponent
becomes a quantity one may choose rather than merely observe. If a construction admits a family of
admissible maps sending the object of interest into $(-1,1)$, the map with the smallest endpoint
exponent yields the numerically most tractable Gram family, at a rate that
Proposition \ref{prop:kappa} makes explicit. The Cayley map \eqref{eq:cayley_map} is canonical for
Li's criterion but is not forced by conditioning considerations, and $\beta=2$ is its price.

We note without proof, since it requires the kernel formalism not developed here, that the same
exponent also fixes the critical Schatten exponent $p_{c}=1/\beta$ of the associated integral
kernel, so that $\beta=2$ gives $p_{c}=1/2$. Whether a kernel lies in a given trace ideal decides
which analytic constructions --- Fredholm determinants, regularized traces --- are licensed at all,
so the exponent has consequences beyond magnitude estimates. We record this as an indication and
develop it elsewhere.

\subsection{Why the extra $\log N$: discreteness, not size}

It is worth isolating what the factor $\log N$ is measuring, because it is the feature that
separates this regime from the classical one. For a measure supported on a compact set $E$ of
positive logarithmic capacity, Szeg\H o--Widom theory gives $h_{n}^{1/n}\to\operatorname{cap}(E)^{2}$,
so $\log D_{N}=\sum_{n<N}\log h_{n}\sim N^{2}\log\operatorname{cap}(E)$: order $N^{2}$, with no
logarithm. The Hilbert matrix, the moment matrix of Lebesgue measure on $[0,1]$, is the standard
instance.

The supports considered here are countable and therefore of capacity zero
\cite{StahlTotik1992,Totik2005Orthogonal}, so $h_{n}^{1/n}\to0$ and the $N^{2}$ law degenerates. The
content of Theorem \ref{thm:endpoint} is that the degeneration is not arbitrary: it is exactly one
factor of $\log N$, with a coefficient that is precisely the rate at which the atoms approach their
accumulation point. In the language of numerical analysis, moment matrices on discrete supports
with an accumulation point are worse conditioned than their continuous counterparts by a factor
$\log N$ in the exponent, and $\beta$ says by how much. This statement is independent of $\zeta$ and
of number theory, and is the form in which the result is likely to be useful outside the present
context.

\subsection{What the law does not do}

For completeness, and because the surrounding literature invites the confusion: nothing here bears
on the Riemann hypothesis. The theorem controls the \emph{magnitude} of determinants; Li's criterion
and Weil positivity are statements about \emph{signs}, and a magnitude asymptotic is compatible with
either sign behaviour. Nothing here proves the irrationality of any zeta value; Corollary
\ref{cor:diophantine} improves a constant in a conditional statement and no more. And nothing here
reaches the $N^{2}$ order at which, by \S\ref{ss:nogo}, the interesting information would have to
live.

\begin{thebibliography}{99}

\bibitem{BombieriLagarias1999}
E. Bombieri and J. C. Lagarias,
\emph{Complements to Li's criterion for the Riemann hypothesis},
J. Number Theory \textbf{77} (1999), no. 2, 274--287.
DOI: 10.1006/jnth.1999.2392.

\bibitem{BergSzwarc2014}
C. Berg and R. Szwarc,
\emph{On the order of indeterminate moment problems},
Adv. Math. \textbf{250} (2014), 105--143.
DOI: 10.1016/j.aim.2013.09.015. arXiv:1310.0247.

\bibitem{BergPedersen1994}
C. Berg and H. L. Pedersen,
\emph{On the order and type of the entire functions associated with an indeterminate Hamburger
moment problem},
Ark. Mat. \textbf{32} (1994), no. 1, 1--11.
DOI: 10.1007/BF02559521.

\bibitem{Chihara1978}
T. S. Chihara,
\emph{An Introduction to Orthogonal Polynomials},
Gordon and Breach, New York, 1978.

\bibitem{HaynesZudilin2015}
A. Haynes and W. Zudilin,
\emph{Hankel determinants of zeta values},
SIGMA \textbf{11} (2015), Paper 101, 5 pp.
DOI: 10.3842/SIGMA.2015.101. arXiv:1510.01901.

\bibitem{Johansson2001}
K. Johansson,
\emph{Discrete orthogonal polynomial ensembles and the Plancherel measure},
Ann. of Math. (2) \textbf{153} (2001), no. 1, 259--296.
DOI: 10.2307/2661375.

\bibitem{Lagarias2007}
J. C. Lagarias,
\emph{Li coefficients for automorphic $L$-functions},
Ann. Inst. Fourier (Grenoble) \textbf{57} (2007), no. 5, 1689--1740.
DOI: 10.5802/aif.2311.

\bibitem{Li1997}
X.-J. Li,
\emph{The positivity of a sequence of numbers and the Riemann hypothesis},
J. Number Theory \textbf{65} (1997), no. 2, 325--333.
DOI: 10.1006/jnth.1997.2137.

\bibitem{Monien2009}
H. Monien,
\emph{Hankel determinants of Dirichlet series},
arXiv:0901.1883 [math.NT] (2009).

\bibitem{StahlTotik1992}
H. Stahl and V. Totik,
\emph{General Orthogonal Polynomials},
Encyclopedia of Mathematics and its Applications, vol. 43,
Cambridge University Press, Cambridge, 1992.
DOI: 10.1017/CBO9780511759420.

\bibitem{Totik2005Orthogonal}
V. Totik,
\emph{Orthogonal polynomials},
Surveys in Approximation Theory \textbf{1} (2005), 70--125.

\bibitem{WuLinXuZhao2016}
X.-B. Wu, Y. Lin, S.-X. Xu and Y.-Q. Zhao,
\emph{Uniform asymptotics for discrete orthogonal polynomials on infinite nodes with an
accumulation point},
Anal. Appl. (Singap.) \textbf{14} (2016), no. 5, 705--737.
DOI: 10.1142/S0219530515500177.

\end{thebibliography}

\end{document}
